\documentclass{dlproblemset}
\usepackage{booktabs}
\usepackage{hyperref}
\usepackage{tikz}
\usetikzlibrary{arrows.meta,decorations.pathreplacing,positioning}
\usepackage{emoji}
\tikzset{every picture/.style={line width=0.75pt}}

\definecolor{CellBlue}{HTML}{BBDEFB}
\definecolor{CellGreen}{HTML}{C8E6C9}
\definecolor{CellYellow}{HTML}{FFF9C4}
\definecolor{CellRed}{HTML}{FFCDD2}
\definecolor{LightGray}{HTML}{F5F5F5}
\definecolor{DarkBlue}{HTML}{1565C0}
\definecolor{DarkGreen}{HTML}{1B5E20}
\definecolor{DarkRed}{HTML}{C62828}

\title{Atenção e Transformers}
\subtitle{Modelos Encoder-Only}

\begin{document}

\paragraph{Instruções:} Cada questão possui exatamente uma alternativa correta.

% ============================================================
\section*{Questões}
% ============================================================

\begin{enumerate}[I)]

% ------------------------------------------------------------------
% Q1 — Fácil (NER: esquema BIO)
% ------------------------------------------------------------------
\item Na tarefa de \textit{Named Entity Recognition} (NER) com esquema
      \textbf{BIO}, as categorias são: \textbf{B}egin (início de entidade),
      \textbf{I}nside (continuação) e \textbf{O}utside. Para a frase

\begin{center}
  \texttt{Lucas trabalha na PUCRS em Porto Alegre}
\end{center}

considerando as entidades \textit{Lucas} (pessoa), \textit{PUCRS} (organização)
e \textit{Porto Alegre} (local, dois \textit{tokens}), qual sequência de
\textit{tags} é correta?

\begin{enumerate}[a)]
    \item \texttt{B-PER, O, O, B-ORG, O, B-LOC, I-LOC}
    \item \texttt{B-PER, O, O, B-ORG, O, I-LOC, I-LOC}
    \item \texttt{B-PER, O, O, I-ORG, O, B-LOC, I-LOC}
    \item \texttt{I-PER, O, O, B-ORG, O, B-LOC, I-LOC}
\end{enumerate}
% R: a

% ------------------------------------------------------------------
% Q2 — Médio (BERT: contagem de parâmetros de atenção)
% ------------------------------------------------------------------
\item O BERT-Base tem $L = 12$ camadas, dimensão oculta $d = 768$ e
      $H = 12$ cabeças de atenção com $d_k = d/H = 64$ por cabeça.
      Em uma única camada de \textit{self-attention}, as projeções são:
      \[
        \mathbf{W}^Q,\; \mathbf{W}^K,\; \mathbf{W}^V \in \mathbb{R}^{d \times (H \cdot d_k)},
        \qquad
        \mathbf{W}^O \in \mathbb{R}^{(H \cdot d_k) \times d}.
      \]
      Ignorando vieses, quantos parâmetros treináveis existem nessas
      \textbf{quatro} matrizes em uma única camada?

\begin{enumerate}[a)]
    \item $589\,824$: apenas $\mathbf{W}^Q$.
    \item $1\,769\,472$: $\mathbf{W}^Q, \mathbf{W}^K, \mathbf{W}^V$ sem a projeção de saída $\mathbf{W}^O$.
    \item $2\,359\,296$: as quatro matrizes $\mathbf{W}^Q, \mathbf{W}^K, \mathbf{W}^V, \mathbf{W}^O$.
    \item $4\,718\,592$: dupla contagem ao considerar dimensões de
          entrada e saída separadamente para cada projeção.
\end{enumerate}
% R: c

% ------------------------------------------------------------------
% Q3 — Médio (MLM: regra 80/10/10 e desalinhamento treino-inferência)
% ------------------------------------------------------------------
\item No pré-treino do BERT, a estratégia MLM seleciona aleatoriamente
      $15\%$ dos \textit{tokens}. Um estudante propõe simplificar a regra
      $80/10/10$ substituindo $\mathbf{100\%}$ dos \textit{tokens}
      selecionados pelo \textit{token} especial \texttt{[MASK]} (em vez
      de $80\%$ por \texttt{[MASK]}, $10\%$ por palavra aleatória e
      $10\%$ mantidos inalterados). Qual é a principal desvantagem
      dessa simplificação?

\begin{enumerate}[a)]
    \item O modelo nunca aprende a representar \textit{tokens} reais
          durante o pré-treino; ao fazer \textit{fine-tuning}, jamais
          encontra \texttt{[MASK]}, criando um desalinhamento entre
          as duas fases.
    \item A perda MLM fica indefinida porque \texttt{[MASK]} não faz
          parte do vocabulário padrão do modelo.
    \item Os gradientes da posição \texttt{[MASK]} se anulam quando
          todas as posições selecionadas recebem o mesmo \textit{token}.
    \item O modelo aprende a ignorar a posição dos \textit{tokens}, pois
          \texttt{[MASK]} é sempre o mesmo símbolo, independentemente
          da posição.
\end{enumerate}
% R: a

% ------------------------------------------------------------------
% Q4 — Médio (BERT: segment embeddings e tarefas par-de-sentença)
% ------------------------------------------------------------------
\item A representação de entrada do BERT é $\mathbf{x}_i = \mathbf{e}_{\text{token}}(w_i)
      + \mathbf{e}_{\text{segmento}}(s_i) + \mathbf{e}_{\text{posição}}(i)$, onde
      $\mathbf{e}_{\text{segmento}}$ assume o valor $\mathbf{e}_A$ para tokens da primeira
      sentença e $\mathbf{e}_B$ para tokens da segunda. Um pesquisador remove
      o componente $\mathbf{e}_{\text{segmento}}$ para simplificar o modelo.
      Em qual das seguintes tarefas essa remoção causaria a
      \textbf{maior degradação} de desempenho?

\begin{enumerate}[a)]
    \item \textit{POS tagging} em uma única sentença.
    \item Classificação de sentimento em uma única sentença (SST-2).
    \item \textit{Extractive QA} (SQuAD): identificar o trecho que responde
          a uma pergunta dado um parágrafo de contexto.
    \item \textit{Named Entity Recognition} (NER) em uma única sentença.
\end{enumerate}
% R: c

% ------------------------------------------------------------------
% Q5 — Médio (ELECTRA vs. BERT: eficiência de treino)
% ------------------------------------------------------------------
\item O ELECTRA (Clark et al., 2020) substitui o objetivo MLM do BERT
      pelo \textit{Replaced Token Detection} (RTD): um gerador pequeno
      cria \textit{tokens} substitutos para $\approx 15\%$ das posições,
      e o discriminador ELECTRA deve classificar \textbf{cada posição}
      da sequência como ``real'' ou ``substituído''. Comparado ao BERT
      com MLM, qual é a principal vantagem do ELECTRA em eficiência
      de treinamento?

\begin{enumerate}[a)]
    \item O ELECTRA não precisa de uma camada de vocabulário na saída
          (custo $d \times |V|$), pois a classificação RTD é binária.
    \item O ELECTRA gera sinal de gradiente em \textbf{todas} as posições
          da sequência (não apenas nos $\approx 15\%$ mascarados do
          MLM), obtendo desempenho equivalente com $\approx 25\%$
          do \textit{compute}.
    \item O ELECTRA não precisa de \textit{fine-tuning}, pois o RTD
          já otimiza diretamente tarefas discriminativas como NER.
    \item O ELECTRA usa \textit{batches} menores que o BERT por precisar
          do gerador auxiliar, reduzindo o consumo de memória e
          acelerando o treino por passo.
\end{enumerate}
% R: b

% ------------------------------------------------------------------
% Q6 — Médio (CLS vanilla falha em STS: anisotropia)
% ------------------------------------------------------------------
\item Um engenheiro extrai $\mathbf{h}_{[\texttt{CLS}]}^{(L)}$ (estado
      oculto final do \texttt{[CLS]}) de um BERT pré-treinado
      \textbf{sem fine-tuning} para medir a similaridade semântica entre
      sentenças via cosseno. No experimento de Reimers \& Gurevych (2019),
      esse método obtém $20.29$ de correlação de Spearman no
      STS-Benchmark, abaixo do GloVe ($58.02$). Qual raciocínio
      explica corretamente esse fenômeno?

\begin{enumerate}[a)]
    \item O \texttt{[CLS]} foi treinado no NSP, um objetivo binário e
          ruidoso que não incentiva a organização do espaço para
          similaridade cosseno; o espaço resultante é anisotrópico
          (vetores concentrados em um cone estreito), tornando as
          distâncias cosseno pouco discriminativas.
    \item O BERT usa \textit{position embeddings} aprendidos, logo
          o \texttt{[CLS]} captura apenas informação posicional e
          perde a semântica do restante da sentença.
    \item O \texttt{[CLS]} usa apenas os \textit{tokens} da última camada,
          enquanto o GloVe usa toda a hierarquia de \textit{n}-grams
          (mais informação leva ao melhor desempenho).
    \item O BERT foi pré-treinado apenas em inglês, e o STS-Benchmark
          contém sentenças em múltiplos idiomas, favorecendo o GloVe
          multilíngue.
\end{enumerate}
% R: a

% ------------------------------------------------------------------
% Q7 — Médio (mean pooling mascarado: cálculo com token [PAD])
% ------------------------------------------------------------------
\item No SBERT, o \textit{mean pooling} mascarado calcula o vetor de
      sentença excluindo posições \texttt{[PAD]}:
      \[
        \mathbf{u} = \frac{\sum_{i=1}^{L} m_i\,\mathbf{h}_i}{\sum_{i=1}^{L} m_i},
        \qquad m_i \in \{0,1\}.
      \]
      Uma sentença tem $L = 3$ tokens com os seguintes \textit{embeddings}
      $\mathbf{h}_i \in \mathbb{R}^2$ e máscara de atenção:

\begin{center}
\begin{tabular}{lcc}
  \toprule
  Token & $\mathbf{h}_i$ & $m_i$ \\
  \midrule
  $w_1$ & $[2,\;0]$ & $1$ \\
  $w_2$ & $[0,\;2]$ & $1$ \\
  $w_3$ (\texttt{[PAD]}) & $[1,\;3]$ & $0$ \\
  \bottomrule
\end{tabular}
\end{center}

Qual é o vetor de sentença $\mathbf{u}$ produzido pelo \textit{mean pooling} mascarado?

\begin{enumerate}[a)]
    \item $[1,\;1]$:
          média dos tokens reais $w_1$ e $w_2$, excluindo o \texttt{[PAD]}.
    \item $\bigl[1,\;\tfrac{5}{3}\bigr] \approx [1.00,\;1.67]$:
          inclui $\mathbf{h}_3$ do \texttt{[PAD]} no numerador e usa $L = 3$
          no denominador.
    \item $[2,\;2]$:
          soma correta dos tokens mascarados, sem dividir por $\sum m_i$.
    \item $\bigl[\tfrac{2}{3},\;\tfrac{2}{3}\bigr] \approx [0.67,\;0.67]$:
          exclui o \texttt{[PAD]} do numerador, mas divide por $L = 3$
          em vez de $\sum m_i = 2$.
\end{enumerate}
% R: a

% ------------------------------------------------------------------
% Q8 — Médio (bi-encoder vs. cross-encoder: latência de consulta)
% ------------------------------------------------------------------
\item Um sistema de busca semântica precisa classificar $n = 100\,000$
      documentos em relação a uma consulta. Um \textbf{cross-encoder}
      requer $1$ \textit{forward pass} por par (consulta, documento);
      um \textbf{bi-encoder} pré-computa os documentos \textit{offline} e,
      no momento da consulta, realiza apenas $1$ \textit{forward pass}
      (para a consulta) seguido de busca vetorial. Cada \textit{forward
      pass} demora $t = 10\text{ ms}$. Qual é a redução de latência no
      momento da consulta ao usar o bi-encoder?

\begin{enumerate}[a)]
    \item $\approx 100\times$: o bi-encoder ainda precisa de $1$
          \textit{forward pass} por documento durante a consulta.
    \item $\approx 1\,000\times$: eliminando o processamento par-a-par
          para os primeiros $1\,000$ documentos filtrados.
    \item $\approx 100\,000\times$: a consulta requer apenas $10\text{ ms}$
          (um único \textit{forward pass}) em vez de
          $100\,000 \times 10\text{ ms} = 1\,000\text{ s}$.
    \item Equivalente; a busca vetorial requer um produto escalar por
          documento, igualando o custo do \textit{forward pass}.
\end{enumerate}
% R: c

% ------------------------------------------------------------------
% Q9 — Médio (InfoNCE: efeito do tamanho de batch)
% ------------------------------------------------------------------
\item Na perda \textit{InfoNCE} com \textit{in-batch negatives}, um
      \textit{batch} de $N = 4$ pares $(x_i, x_i^+)$ fornece $N - 1 = 3$
      negativos por âncora. Um engenheiro aumenta o \textit{batch}
      para $N = 1\,024$. Qual é o efeito esperado desse aumento?

\begin{enumerate}[a)]
    \item O treino fica mais lento por \textit{token}, mas os gradientes
          são mais informativos: cada âncora deve distinguir seu positivo
          de $1\,023$ competidores, melhorando a qualidade dos
          \textit{embeddings}.
    \item O treino produz resultado equivalente ao de $N = 4$: a perda
          InfoNCE é normalizada pelo número de negativos e compensa
          automaticamente o aumento.
    \item O treino diverge com $N$ grande porque o denominador do
          \textit{softmax} fica muito grande, zerando os gradientes.
    \item O benefício de $N$ grande é apenas computacional (melhor
          utilização de GPU), sem impacto na qualidade final dos
          \textit{embeddings}.
\end{enumerate}
% R: a

% ------------------------------------------------------------------
% Q10 — Difícil (triplet loss: cálculo numérico)
% ------------------------------------------------------------------
\item \emoji{skull-and-crossbones} Considere três \textit{embeddings}
      (já normalizados) e a \textit{triplet loss} com margem $\alpha = 0.5$:
      \[
        \mathbf{f}(x) = [1.0,\;0.0], \quad
        \mathbf{f}(x^+) = [0.6,\;0.8], \quad
        \mathbf{f}(x^-) = [-0.6,\;0.8],
      \]
      \[
        \mathcal{L}_{\text{triplet}} = \max\!\bigl(0,\;
          \|\mathbf{f}(x)-\mathbf{f}(x^+)\|_2^2 - \|\mathbf{f}(x)-\mathbf{f}(x^-)\|_2^2 + \alpha\bigr).
      \]
      Qual é o valor de $\mathcal{L}_{\text{triplet}}$?

\begin{enumerate}[a)]
    \item $0$ (a margem já é satisfeita: o positivo está suficientemente
          mais próximo que o negativo).
    \item $0.50$ (a perda iguala exatamente a margem $\alpha$ quando
          as duas distâncias são iguais).
    \item $1.90$ (diferença das distâncias quadráticas,
          sem aplicar o $\max(0, \cdot)$).
    \item $-1.90$ (resultado antes do truncamento por zero).
\end{enumerate}
% R: a

% ------------------------------------------------------------------
% Q11 — Difícil (MLM loss: cálculo numérico)
% ------------------------------------------------------------------
\item \emoji{skull-and-crossbones} O BERT pré-treina com a perda MLM:
      \[
        \mathcal{L}_{\text{MLM}}
          = -\frac{1}{|M|}\sum_{i \in M} \log p(w_i \mid \mathbf{w}_{\setminus M}).
      \]
      Em um passo de treino, $|M| = 3$ posições são mascaradas. O modelo
      atribui as seguintes probabilidades ao \textit{token} correto em cada
      posição: $p_1 = 0.50$,\ $p_2 = 0.25$,\ $p_3 = 0.125$.
      Qual é o valor de $\mathcal{L}_{\text{MLM}}$?

\begin{enumerate}[a)]
    \item $\approx 0.693$
    \item $\approx 1.386$
    \item $\approx 2.079$
    \item $\approx 4.159$
\end{enumerate}
% R: b

\end{enumerate}


% ============================================================
\section*{Soluções}
% ============================================================

% --- Solução Q1 ---
\subsection*{I — Resposta: (a)}
O esquema BIO exige que o \textbf{primeiro token} de qualquer entidade receba
a \textit{tag} \textbf{B-} (Begin), e os tokens subsequentes da mesma entidade
recebam \textbf{I-} (Inside). Tokens fora de entidades recebem \textbf{O}.
Na frase:
\begin{itemize}
  \item ``Lucas'' $\to$ \texttt{B-PER} (inicia a entidade pessoa).
  \item ``PUCRS'' $\to$ \texttt{B-ORG} (inicia a entidade organização).
  \item ``Porto'' $\to$ \texttt{B-LOC} (inicia a entidade local de dois tokens).
  \item ``Alegre'' $\to$ \texttt{I-LOC} (continua a mesma entidade local).
\end{itemize}
A opção~(b) usa \texttt{I-LOC} para ``Porto'' (erro: nunca se inicia uma
entidade com \texttt{I-}).
A opção~(c) usa \texttt{I-ORG} para ``PUCRS'' (mesmo erro para organização).
A opção~(d) usa \texttt{I-PER} para ``Lucas'' (mesmo erro para pessoa).

% --- Solução Q2 ---
\subsection*{II — Resposta: (c)}
Como $H \cdot d_k = 12 \times 64 = 768 = d$, todas as quatro matrizes são
$\mathbb{R}^{768 \times 768}$. O número de parâmetros em cada uma é
$768 \times 768 = 589\,824$, e no total:
\[
  4 \times 589\,824 = 2\,359\,296.
\]
A opção~(a) conta apenas $\mathbf{W}^Q$.
A opção~(b) omite a projeção de saída $\mathbf{W}^O$, que transforma a concatenação das
cabeças de volta à dimensão $d$ e é tão grande quanto as demais.
A opção~(d) duplica o resultado correto (erro de dupla contagem).

% --- Solução Q3 ---
\subsection*{III — Resposta: (a)}
O \texttt{[MASK]} é um token \textbf{exclusivo do pré-treino}: ele nunca aparece
na fase de \textit{fine-tuning} nem na inferência. Se todos os tokens selecionados
fossem substituídos por \texttt{[MASK]}, o modelo \textbf{nunca aprenderia a
representar tokens reais}; ao encontrá-los no \textit{fine-tuning}, suas
representações seriam subótimas. A regra $80/10/10$ mitiga esse desalinhamento
ao forçar o modelo a lidar com tokens reais (intactos ou aleatórios) mesmo durante
o pré-treino.
A opção~(b) é falsa: \texttt{[MASK]} é um token do vocabulário do BERT.
A opção~(c) é falsa: posições com o mesmo \textit{token} ainda produzem gradiente
distinto, pois os estados ocultos dependem do contexto.
A opção~(d) confunde mascaramento com ignorância de posição; a \textit{position
embedding} é somada antes de qualquer mascaramento.

% --- Solução Q4 ---
\subsection*{IV — Resposta: (c)}
No formato SQuAD, a entrada é \texttt{[CLS] pergunta [SEP] contexto [SEP]}, com
os tokens da pergunta recebendo $\mathbf{e}_A$ e os do contexto recebendo $\mathbf{e}_B$. Sem
\textit{segment embeddings}, o modelo \textbf{não distingue} quais tokens são a
pergunta e quais são o contexto, degradando drasticamente a tarefa de
identificação de \textit{spans}.
As tarefas (a), (b) e (d) operam em uma única sentença, portanto todos os tokens
compartilhariam $\mathbf{e}_A$: a remoção de $\mathbf{e}_{\text{segmento}}$ não altera
o comportamento nessas tarefas.

% --- Solução Q5 ---
\subsection*{V — Resposta: (b)}
No MLM do BERT, o gradiente só flui pelas $\approx 15\%$ das posições
mascaradas. No ELECTRA, o discriminador classifica \textbf{todas} as $L$ posições
(``real'' ou ``substituído''), gerando gradiente \textbf{denso} a cada passo.
Esse sinal mais rico permite que o ELECTRA alcance o mesmo desempenho do BERT
com $\approx 25\%$ do \textit{compute} de pré-treino.
A opção~(a) é parcialmente verdadeira (a cabeça de saída é binária), mas a
camada de vocabulário do gerador ainda existe; a vantagem principal é o sinal denso.
A opção~(c) é falsa: o ELECTRA ainda requer \textit{fine-tuning}.
A opção~(d) inverte o raciocínio: o gerador aumenta o custo por passo, mas o
ganho em eficiência vem do sinal mais informativo.

% --- Solução Q6 ---
\subsection*{VI — Resposta: (a)}
O \texttt{[CLS]} é treinado no NSP (classificação binária ``é ou não é a
próxima sentença''), sem nenhum incentivo para estruturar o espaço de
\textit{embeddings} de forma que \textbf{frases similares fiquem próximas}.
O resultado é um espaço \textbf{anisotrópico}: todos os vetores se concentram
em um cone estreito do hiperespaço, tornando as distâncias cosseno pouco
discriminativas (frases totalmente diferentes podem ter cosseno $> 0.9$).
O GloVe, embora estático, captura co-ocorrências e produz um espaço
mais isotrópico, o que o torna superior para STS \textit{vanilla}.
A opção~(b) confunde \textit{position embeddings} com o objetivo do
\texttt{[CLS]}.
A opção~(c) inventa um mecanismo inexistente.
A opção~(d) é falsa: o BERT original foi pré-treinado em inglês, assim como o
GloVe mais comum, e o STS-Benchmark é predominantemente em inglês.

% --- Solução Q7 ---
\subsection*{VII — Resposta: (a)}
A máscara zera a contribuição de $w_3$; o numerador é
$m_1\mathbf{h}_1 + m_2\mathbf{h}_2 + m_3\mathbf{h}_3
 = [2,\,0] + [0,\,2] + \mathbf{0} = [2,\,2]$
e o denominador é $\sum m_i = 1 + 1 + 0 = 2$:
\[
  \mathbf{u} = \frac{[2,\,2]}{2} = [1,\,1].
\]
A opção~(b) inclui $\mathbf{h}_3 = [1,\,3]$ do \texttt{[PAD]} no numerador e usa $L = 3$:
$([2,0]+[0,2]+[1,3])/3 = [3,5]/3 = [1,\,\tfrac{5}{3}]$: resultado distorcido pela posição de \textit{padding}.
A opção~(c) aplica corretamente a máscara ao numerador, mas omite a divisão por $\sum m_i$,
produzindo a soma bruta $[2,\,2]$.
A opção~(d) exclui o \texttt{[PAD]} do numerador (correto), mas divide por $L = 3$ em vez
de $\sum m_i = 2$, gerando $[2,\,2]/3 = [\tfrac{2}{3},\,\tfrac{2}{3}]$: erro de normalização.

% --- Solução Q8 ---
\subsection*{VIII — Resposta: (c)}
O cross-encoder requer $1$ \textit{forward pass} por par, logo $n = 100\,000$
passes no momento da consulta:
\[
  100\,000 \times 10\text{ ms} = 1\,000\,000\text{ ms} = 1\,000\text{ s}.
\]
O bi-encoder pré-computa todos os documentos \textit{offline}. Na consulta, basta
$1$ \textit{forward pass} para a \textit{query} ($10\text{ ms}$) seguido de busca
vetorial (produto escalar em hardware SIMD, $\ll 1\text{ ms}$ para $10^5$ vetores).
Redução de latência:
\[
  \frac{1\,000\,000\text{ ms}}{10\text{ ms}} = 100\,000\times.
\]
A opção~(a) erra ao supor que o bi-encoder ainda faz um \textit{forward pass} por
documento em tempo de consulta (documentos estão pré-computados).
A opção~(b) chega a $1\,000\times$ por confundir $n = 100\,000$ com $n = 1\,000$.
A opção~(d) subestima a diferença de velocidade entre um produto escalar e um
\textit{forward pass} completo por um transformer (fator de $\sim 10^4$--$10^6$).

% --- Solução Q9 ---
\subsection*{IX — Resposta: (a)}
Na perda InfoNCE, para cada âncora $x_i$ o positivo $x_i^+$ deve ser identificado
entre todos os $x_j^+$ do \textit{batch}:
\[
  \mathcal{L}_{\text{NCE}} = -\frac{1}{N}\sum_{i=1}^{N} \log
    \frac{e^{\text{sim}(\mathbf{f}(x_i),\,\mathbf{f}(x_i^+))/\tau}}
         {\sum_{j=1}^{N} e^{\text{sim}(\mathbf{f}(x_i),\,\mathbf{f}(x_j^+))/\tau}}.
\]
Com $N = 1\,024$, cada âncora compete com $1\,023$ negativos simultâneos, gerando
gradiente mais rico e discriminativo. Empiricamente, o desempenho escala com
$\log N$, motivando \textit{batches} de $1\,024$--$16\,384$ em modelos SOTA.
A opção~(b) erra: a normalização pelo $N$ no denominador da média \textbf{não} é
equivalente a neutralizar o efeito de mais negativos: o denominador do
\textit{softmax} cresce com $N$, tornando o problema de classificação mais difícil.
A opção~(c) erra: o \textit{softmax} com denominador grande concentra gradiente
nos negativos mais difíceis, não o zera.
A opção~(d) erra: mais negativos melhoram a qualidade, não apenas a eficiência.

% --- Solução Q10 ---
\subsection*{X — Resposta: (a)}
\textbf{Cálculo das distâncias quadráticas:}
\begin{align*}
  \|\mathbf{f}(x) - \mathbf{f}(x^+)\|_2^2
    &= (1.0 - 0.6)^2 + (0.0 - 0.8)^2
     = 0.16 + 0.64 = 0.80, \\
  \|\mathbf{f}(x) - \mathbf{f}(x^-)\|_2^2
    &= (1.0 - (-0.6))^2 + (0.0 - 0.8)^2
     = 2.56 + 0.64 = 3.20.
\end{align*}
\textbf{Aplicando a fórmula:}
\[
  \mathcal{L}_{\text{triplet}}
    = \max(0,\; 0.80 - 3.20 + 0.50)
    = \max(0,\; -1.90)
    = 0.
\]
O positivo ($d^2 = 0.80$) já está $2.40$ unidades mais próximo que o negativo
($d^2 = 3.20$), excedendo a margem $\alpha = 0.50$: o par não viola a
restrição e não gera gradiente.
A opção~(b) ocorreria se as duas distâncias fossem iguais ($d^+ = d^-$); nesse
caso $\mathcal{L} = \alpha = 0.50$.
A opção~(c) é o argumento antes do $\max(0, \cdot)$.
A opção~(d) é o mesmo sem o valor absoluto.

% --- Solução Q11 ---
\subsection*{XI — Resposta: (b)}
As probabilidades dos tokens corretos são $p_1 = \tfrac{1}{2}$,
$p_2 = \tfrac{1}{4}$, $p_3 = \tfrac{1}{8}$. Aplicando a fórmula com $|M| = 3$:
\begin{align*}
  \mathcal{L}_{\text{MLM}}
    &= -\frac{1}{3}\Bigl(\ln\tfrac{1}{2} + \ln\tfrac{1}{4} + \ln\tfrac{1}{8}\Bigr) \\
    &= -\frac{1}{3}\Bigl(-\ln 2 - 2\ln 2 - 3\ln 2\Bigr) \\
    &= \frac{6\ln 2}{3} = 2\ln 2 \approx 1.386.
\end{align*}
\textbf{Análise das alternativas incorretas:}
\begin{itemize}
  \item \emph{(a) $\approx 0.693$}: considera apenas $-\ln(0.50) = \ln 2$,
        ou seja, somente a primeira posição, sem a média.
  \item \emph{(c) $\approx 2.079$}: usa $|M| = 2$ no denominador ($6\ln 2 / 2 = 3\ln 2$),
        como se apenas duas posições fossem mascaradas.
  \item \emph{(d) $\approx 4.159$}: omite a normalização $1/|M|$ por completo
        ($-(\ln\tfrac{1}{2}+\ln\tfrac{1}{4}+\ln\tfrac{1}{8}) = 6\ln 2$).
\end{itemize}

\end{document}
